\section{Gaussian 与多元 Gaussian}
\label{sec:05-Probability/06-Common-Distributions/05}

训练完一个回归模型，你把残差画成直方图，得到一个还算周正的钟形。你用的损失函数是平方误差——这等价于假设残差服从 Gaussian，只是没人在代码里写出来。

问题是：为什么这是默认？换成绝对值损失就是假设 Laplace 分布，换成 Huber 就是别的东西。默认背后应该有理由，而 Gaussian 恰好有两个，彼此独立，各自也各有失效的条件。

第一个理由是加法。如果你观测到的偏差是许多互不相干的小因素叠加出来的——测量噪声、采样波动、上游服务的抖动、无数个没进模型的微小变量——那么不管每个因素自己长什么样，它们的和都会向同一个形状靠拢。

\begin{center}
\begin{tikzpicture}[x=1cm,y=1cm,font=\footnotesize]
  \foreach \dx/\lab in {0/{一个}, 5.0/{两个之和}, 10.0/{三个之和}}{
    \draw[black!60] ({\dx-2.2},0) -- ({\dx+2.2},0);
    \node[black!70] at (\dx,-0.5) {\lab};}
  \draw[blue!55!black,thick] (-0.7,0) -- (-0.7,2) -- (0.7,2) -- (0.7,0);
  \begin{scope}[shift={(5.0,0)}]
    \draw[blue!55!black,thick] (-1.4,0) -- (0,2) -- (1.4,0);
  \end{scope}
  \begin{scope}[shift={(10.0,0)}]
    \draw[black!45,dashed,thick,domain=-1.85:1.85,samples=70]
      plot ({1.4*\x},{2*0.79788*exp(-2*\x*\x)});
    \draw[blue!55!black,thick,domain=-1.5:-0.5,samples=25]
      plot ({1.4*\x},{2*0.5*(1.5+\x)*(1.5+\x)});
    \draw[blue!55!black,thick,domain=-0.5:0.5,samples=25]
      plot ({1.4*\x},{2*(0.75-\x*\x)});
    \draw[blue!55!black,thick,domain=0.5:1.5,samples=25]
      plot ({1.4*\x},{2*0.5*(1.5-\x)*(1.5-\x)});
    \node[black!55] at (2.0,1.9) {Gaussian};
  \end{scope}
\end{tikzpicture}

\emph{起点是一个方方正正的均匀分布，和它自己相加一次变成三角形，再加一次就已经贴着虚线画的 Gaussian 了。加法把棱角磨掉的速度快得出奇——第三步的最大偏差不到峰值的百分之七。}
\end{center}

第二个理由和加法无关。在只知道均值和方差、别的一概不知的情况下，Gaussian 是熵最大的分布，也就是在这两条约束下承诺最少的那个（推导见\hyperref[sec:07-Information-Theory/01-Information-and-Entropy/04]{``最大熵原理：从约束到分布''}）。选它不是因为你相信数据是钟形的，而是因为你不想在没有证据的地方额外假定结构。

两个理由的失效条件不同，这一点值得先记下来。第一个理由要求那些小因素的方差有限且没有哪一个占压倒优势——这在重尾场景里直接崩掉，本章第七节会详细算这笔账。第二个理由要求你的约束确实只有均值和方差；如果你还知道数据非负，最大熵分布就变成指数型的，知道数据有界则变成均匀型的。默认值只在默认的前提下才是好的。

\subsection{一维：一个位置参数和一个尺度参数}

\(X\sim\mathcal N(\mu,\sigma^2)\) 的密度是

\[
f(x)=\frac{1}{\sqrt{2\pi\sigma^2}}\exp\Bigl[-\frac{(x-\mu)^2}{2\sigma^2}\Bigr],
\]

\(\E[X]=\mu\)，\(\Var(X)=\sigma^2\)，密度关于 \(\mu\) 对称，均值、中位数、众数三点重合。标准化 \(Z=(X-\mu)/\sigma\) 把任何 Gaussian 拉回 \(\mathcal N(0,1)\)，于是全部区间概率都能用一张标准正态表回答。

\(\mu\pm\sigma\)、\(\mu\pm2\sigma\)、\(\mu\pm3\sigma\) 分别覆盖约 \(68\%\)、\(95\%\)、\(99.7\%\)。这三个数是拿来口算的，不是拿来做尾部推断的——\(\mu\pm2\sigma\) 的精确值是 \(95.45\%\)，那 \(0.45\%\) 的差别在构造置信区间时会现形。更要紧的是它们只在数据确实近似 Gaussian 时才有意义，而这恰恰是需要检验而不是假定的。

线性组合的封闭性是 Gaussian 最省心的地方：独立 Gaussian 的任意线性组合仍是 Gaussian，均值按系数加权相加，方差按系数平方加权相加。差也是和——\(X_1-X_2\) 的方差是 \(\sigma_1^2+\sigma_2^2\) 而不是相减，两个不确定性来源无论你关心和还是差，总的不确定性都在累积。

推广到向量时，正确的定义不是``每个分量都是 Gaussian''，而是：\(X\in\R^d\) 称为联合 Gaussian，当且仅当对任意 \(a\in\R^d\)，标量 \(a^{\trans}X\) 都服从一维 Gaussian（允许方差为零）。这个说法绕开了密度，即使协方差矩阵奇异、密度不存在，概念仍然成立。

\subsection{协方差矩阵决定等高线的朝向与粗细}

协方差矩阵 \(\Sigma\) 正定时密度写得出来：

\[
f(x)=\frac{1}{(2\pi)^{d/2}\abs{\Sigma}^{1/2}}
\exp\Bigl[-\tfrac12(x-\mu)^{\trans}\Sigma^{-1}(x-\mu)\Bigr].
\]

密度只通过指数里那个二次型依赖 \(x\)，所以等密度面就是二次型的等值面，也就是以 \(\mu\) 为中心的椭球。二次型 \(\Delta^2=(x-\mu)^{\trans}\Sigma^{-1}(x-\mu)\) 叫平方 Mahalanobis 距离：\(\Sigma=I\) 时它退化成普通的欧氏距离平方，一般情形下它按各方向的散布把距离重新标定，使``等距离''与``等密度''重合。

椭球的几何完全由 \(\Sigma\) 的谱分解读出。\(\Sigma\) 对称，所以它的特征向量正交、特征值实且非负（这一步的完整依据见\hyperref[sec:03-Linear-Algebra/06-Symmetric-and-Positive-Definite-Matrices/01]{``对称矩阵与谱定理：什么时候特征方向一定正交''}）。特征向量给出椭球主轴的方向，特征值的平方根给出各主轴的长度，行列式则是全部特征值之积，衡量数据云占了多大体积。

\begin{center}
\begin{tikzpicture}[x=1cm,y=1cm,font=\footnotesize,>=stealth]
  \draw[->,black!55] (-3.4,0) -- (3.6,0);
  \node[right,black!60] at (3.6,0) {身高偏离};
  \draw[->,black!55] (0,-3.0) -- (0,3.4);
  \node[above,black!60] at (0,3.4) {体重偏离};
  \begin{scope}[rotate=45]
    \draw[blue!50!black,fill=blue!12] (0,0) ellipse (1.73 and 0.73);
    \draw[blue!50!black,dashed] (0,0) ellipse (2.60 and 1.10);
    \draw[->,red!65!black,thick] (0,0) -- (1.73,0);
    \draw[->,red!65!black,thick] (0,0) -- (0,0.73);
  \end{scope}
  \node[red!65!black] at (2.75,2.15) {\(v_1\)，\(\sqrt{\lambda_1}=13.0\)};
  \node[red!65!black] at (-2.05,1.55) {\(v_2\)，\(\sqrt{\lambda_2}=5.5\)};
\end{tikzpicture}

\emph{取 \(\Sigma=\bigl[\begin{smallmatrix}100&70\\70&100\end{smallmatrix}\bigr]\)。特征值是 \(170\) 和 \(30\)，特征向量沿 \(\pm45^\circ\)。长轴指向``整体块头''这个方向，身高体重同增同减；短轴指向``瘦高还是矮胖''，散布只有长轴的四成多。实线是一倍标准差的等高线，虚线是一倍半。协方差矩阵里那个 \(70\) 全部体现为椭圆的倾斜。}
\end{center}

椭球这个画面也解释了两个常用操作。取 \(A=\Sigma^{-1/2}\) 做线性变换，就把椭球拉回单位球，得到 \(\mathcal N(0,I)\)，这叫白化，是主成分分析之类方法的标准预处理。而平方 Mahalanobis 距离本身服从 \(\chi^2_d\)（下一节会说明为什么），于是``离群点''有了一个不依赖坐标尺度的判据：算每个点的 Mahalanobis 距离，和 \(\chi^2_d\) 的高分位数比。一维时它就退回熟悉的``超过几个标准差''。

若 \(\Sigma\) 只是半正定，某些特征值为零，椭球在那些方向上被压扁成一片，全部质量落在一个低维仿射子空间里。此时 \(\R^d\) 上没有密度，但分布本身照样良定义。成分数据（各分量之和恒为一）就是这种情形的常客。

\subsection{条件分布是线性的，而且方差与观测值无关}

把向量分成两块，假定 \(\Sigma_{22}\) 可逆，则

\[
X_1\given X_2=x_2\ \sim\
\mathcal N\bigl(\mu_1+\Sigma_{12}\Sigma_{22}^{-1}(x_2-\mu_2),\;
\Sigma_{11}-\Sigma_{12}\Sigma_{22}^{-1}\Sigma_{21}\bigr).
\]

两件事从这个式子里跳出来。条件均值是 \(x_2\) 的仿射函数，系数矩阵 \(\Sigma_{12}\Sigma_{22}^{-1}\) 扮演回归系数——这就是为什么在 Gaussian 假设下，线性预测已经达到了条件期望的水平，不存在能做得更好的非线性模型。条件协方差里没有 \(x_2\)：无论你观测到什么值，不确定性都减少同样多。这个性质在别的分布里基本不成立，它也是 Kalman 滤波能写成一组固定递推的原因。

拿身高体重的数字走一遍。\(\mu=(170,65)\)，\(\Sigma\) 如上图，相关系数 \(\Corr=0.7\)。得知某人体重 \(75\) 公斤（高出均值 \(10\)），条件均值是 \(170+0.7\times10=177\)，条件方差是 \(100-49=51\)，标准差从 \(10\) 厘米缩到约 \(7.1\)。再看一个反直觉的情形：若得知体重恰好是均值 \(65\)，条件均值仍是 \(170\)，没有任何修正，可条件方差照样从 \(100\) 掉到 \(51\)。观测值不提供修正方向，观测这件事本身仍然消除了不确定性。

顺带补上一条经常被误用的推理。对联合 Gaussian，\(\Sigma_{12}=0\) 时联合密度分解成两个边缘密度之积，不相关等价于独立。这条结论离开``联合''二字就不成立，下面的反例很短：取 \(X\sim\mathcal N(0,1)\)，\(S\) 独立于 \(X\) 且等概率取 \(\pm1\)，令 \(Y=SX\)。由对称性 \(Y\) 也是 \(\mathcal N(0,1)\)，且 \(\E[XY]=\E[S]\E[X^2]=0\)，两者不相关。但 \((X,Y)\) 的全部质量落在 \(y=\pm x\) 两条直线上，给定 \(X=x\) 时 \(Y\) 只能取 \(\pm x\)，是个两点分布。两个 Gaussian 边缘、零相关，联合分布却离二维 Gaussian 十万八千里。协方差矩阵看不穿联合结构，这一点在只报相关系数的分析里很容易被忽略。

\subsection{维度一高，直觉就开始失灵}

标准多元 Gaussian \(\mathcal N(0,I_d)\) 的每个分量都集中在原点附近，可整个向量的范数却集中在 \(\sqrt d\) 附近：\(\norm{X}^2\sim\chi^2_d\)，均值 \(d\)，标准差 \(\sqrt{2d}\)，相对波动 \(\sqrt{2/d}\) 随维度衰减。也就是说，高维 Gaussian 的质量几乎全部住在一层半径约 \(\sqrt d\) 的薄球壳上，而不是原点周围的一团云。

这个现象直接影响两类常见做法。基于欧氏距离的最近邻在高维 Gaussian 数据里会失去分辨力，因为所有点到查询点的距离都差不多；而协方差矩阵的估计在样本量与维度可比时会严重失真，样本特征值向两端散开（详见\hyperref[sec:11-Mathematical-Statistics/09-Random-Matrix-Theory/02]{``Marchenko--Pastur 律给出谱的精确形状''}）。低维画出来的椭圆图景，在高维要谨慎使用。

\subsection{什么时候手上的数据不该用它}

几种迹象出现时值得停下来。数据被硬边界卡住——计数、时长、比例天然非负或有上下限，而 Gaussian 的支撑是整条实轴，尾部会漏到不可能的区域。分布明显偏斜，比如收入或延迟，用均值加方差概括会丢掉最重要的结构。直方图出现不止一个峰，多半意味着混进了不同的子群体，该用混合模型而不是单个 Gaussian。方差随均值变化的异方差数据，与 Gaussian 恒定条件方差的假设直接冲突。极端值出现的频率远超预测——这一条最危险，因为中心区域可能拟合得很好，问题全在尾部。

诊断不必只靠肉眼。Q--Q 图把样本分位数对理论分位数作图，两端向上下翘出对角线说明尾部更重，整体呈弧形说明偏斜，中段出现台阶说明多峰。形式化检验也有，但在大样本下会对毫无实际意义的微小偏离报显著，读结论时要结合偏离的幅度而不只看 \(p\) 值。

出路同样不止一条：对数或 Box--Cox 变换把右偏数据拉回对称，混合 Gaussian 处理多峰，Student \(t\) 给尾部更多容忍，广义线性模型把方差与均值的关系写进模型，实在没有合适参数族时退回核密度估计或直接用分位数。

Gaussian 的位置是基准而不是终点。它让你在线性、可加、对称的世界里迅速得到闭式答案，代价是这三条假设必须真的成立。而下一节会揭示一个更微妙的问题：即使数据确实来自 Gaussian，一旦你用同一批样本同时估计它的均值和方差，你手上那个统计量就已经不是 Gaussian 了。
